\scanpage{7}
\mspage{18}
\diagram{
\begin{tikzcd}
\mathrm{Gr}^{i}(X_{\bar t})_{\underline Q} \arrow[r,"cl"] \arrow[d] & H^{2i}(X_{\bar t},\underline Q_{\ell}(i)) \arrow[d,"\wr"] \\
\mathrm{Gr}^{i}(X_{\bar s})_{\underline Q} \arrow[r,"cl"] & H^{2i}(X_{\bar s},\underline Q_{\ell}(i)).
\end{tikzcd}\qquad (7.13.10)
}

\deleted{\unsure{7.13.11}} \uline{Remarque\deleted{\unsure{s}}} 7.13.11. La commutativité des deux derniers diagrammes \deleted{commutatifs}, qu'on peut exprimer en disant que la spécialisation des cycles algébriques est compatible avec la formation des classes de cohomologie $\ell$-adiques associées, n'est pas très satisfaisante sous la forme énoncée, puisqu'on \deleted{\unsure{a}} a été obligé de remplacer les coefficients $\underline Z_{\ell}(i)$ par $\underline Q_{\ell}(i)$, c'est à dire de négliger la torsion. Pour avoir un énoncé qui ne néglige pas la torsion, \deleted{on est}\inserted{il semble qu'on sera} obligé de travailler avec l'anneau de Chow, puisque c'est sur ce dernier que ce trouve défini naturellement l'homomorphisme $cl$ comme homomorphisme à valeurs dans $H^{2*}(-,\underline Z_{\ell}(*))$, \deleted{\unsure{Mais|Mal-} \unsure{heureusement}} comparer (XIV 4.1). Malheureusement, on n'a pas développé pour le moment une théorie de la spécialisation pour l'anneau de Chow\deleted{\unsure{,}} (cf. 7.14), qui permettrait d'écrire \deleted{un diagramme}\inserted{des carrés} commutatifs analogues à (7.13.9) et (7.13.10) en termes des anneaux de Chow et de la cohomologie à coefficients $\underline Z_{\ell}(*)$. Tout ce qu'on peut dire, \inserted{lorsque $f:X\longrightarrow S$ est lisse et propre} c'est qu'on a un carré commutatif
\diagram{
\begin{tikzcd}
Z^i(X_t) \arrow[r,"cl"] \arrow[d,"\sigma"'] & H^{2i}(X_t,\underline Z_{\ell}(i)) \\
Z^i(X_s) \arrow[r,"cl"] & H^{2i}(X_s,\underline Z_{\ell}(i)), \arrow[u,"\deleted{\unsure{\downarrow}}" description,pos=.15]
\end{tikzcd}\qquad (7.13.12)
}
\deleted{\unsure{et un carré}} analogue à (7.13.9), et un carré correspondant analogue à (7.13.10) qu'on laisse au\deleted{\unsure{x}} lecteur \inserted{le soin} d'écrire, où $Z^i$ désigne le groupe des cycles algébriques (modulo rien du tout) de codimension $i$, \deleted{et} où les flèches horizontales sont les homomorphismes ``classe de cohomologie associée'' définie dans (SGA 5 IV), et où $\sigma$ désigne l'homomorphisme de spécialisation des cycles algébriques, associant à un cycle $Z_t$ le cycle $Z_s$, où $Z$ est

\mspage{6}
